\section{Prewritten answer: decomposition vs.\ trade-off}
% source: cheat-sheet/cheat-sheet.md (Anders's verbose version of the prof's 1-pt solution)

\emph{Q:} Why prefer ``bias--variance \emph{decomposition}'' over ``bias--variance \emph{trade-off}''?

\emph{A (A-level, tight):} ``Decomposition'' is preferred because the identity
$\E[(y_0-\hat f(x_0))^2] = \Bias(\hat f(x_0))^2 + \Var(\hat f(x_0)) + \Var(\varepsilon)$
is an algebraic equality that holds for \emph{every} estimator --- not a constraint forcing one
term up when the other goes down. Two angles make the ``trade-off'' label misleading:
(i)~the bias term is \emph{squared}, so a small absolute bias contributes little to MSE while
variance can shrink a lot (this is exactly why regularisation helps); and
(ii)~changing the model class --- e.g.\ moving into the over-parameterised /
double-descent regime of large NNs --- can lower variance \emph{without} a matching rise
in bias. Stating either angle clearly earns the 1\,P.
